\documentclass[14pt]{extarticle}
\usepackage{fontspec}
\usepackage{polyglossia}
\usepackage{amsmath}
\usepackage{braket}
\usepackage{amsfonts}
\usepackage{amssymb}
\usepackage{pgfplots}
\pgfplotsset{compat=1.18}

% Hebrew setup
% \setdefaultlanguage{hebrew}
% \setotherlanguage{english}
% \newfontfamily\hebrewfont[Script=Hebrew]{Arial}

% Math font setup for proper Greek symbols
\usepackage{unicode-math}
\setmathfont{Latin Modern Math}

% Increase line spacing throughout document
\linespread{1.3}

\usepackage[margin=1in]{geometry}

% Left-aligned equation environment
\newenvironment{lefteqs}{%
    \begin{flushleft}%
    \renewcommand{\arraystretch}{1.3}%
    $\begin{array}{@{}l@{}}%
}{%
    \end{array}$%
    \end{flushleft}%
}

\begin{document}

\begin{center}
    {\LARGE \textbf{Assignment 1 - Quantum Physics 2}}\\[0.5em]
    {\large Shai Roemi - 318986270}\\[0.3em]
\end{center}

\hrule
\vspace{2em}

\textbf{\large 1.}
(Problem 6.23 from course book)

As the questions states, a $\ket{J M}$ coupled state can be represented as an uncoupled state using a Clebsch-Gordan coefficient:

\begin{lefteqs}
    (6.64) \, \ket{J M} = \sum_{m_1, m_2} C^{j_1 j_2 J}_{m_1 m_2 M} \ket{j_1 \, j_2 \, m_1 \, m_2}
\end{lefteqs}

\textbf{(a)} Mutiply from the left by $\bra{j_1 \, j_2 \, m_1 \, m_2}$ to obtain:

\begin{lefteqs}
    \braket{j_1 \, j_2 \, m_1 \, m_2 | J M} = \bra{j_1 \, j_2 \, m_1 \, m_2} \left( \sum_{m_1', m_2'} C^{j_1 j_2 J}_{m_1' m_2' M}  \ket{j_1 \, j_2 \, m_1' \, m_2'}  \right) \\
\end{lefteqs}

The uncoupled states are eigenstates of $\hat{J}_{1z}$ and $\hat{J}_{2z}$ with eigenvalues $m_1 \hbar$ and $m_2 \hbar$ respectively,
where both operators are Hermitian. Eigenvectors of a Hermitian operator that belong to different eigenvalues are orthogonal.

For example, given a Hermitian operator $\hat{A}$ where:
\begin{lefteqs}
    \text{Let } a \neq b \\
    \hat{A} \ket{a} = a \ket{a} \\
    \hat{A} \ket{b} = b \ket{b} \\
    \braket{a | \hat{A} | b} = \bra{a} \left( \hat{A} \ket{b} \right) = \bra{a} b \ket{b} = b \braket{a | b} \\
    \braket{a | \hat{A} | b} = \left( \bra{a} \hat{A} \right) \ket{b} = a^* \braket{a | b} \\
    \braket{a | \hat{A} | b} = a \braket{a | b} \quad\text{(eigenvalues of a Hermitian operator are real)} \\
    (a - b) \braket{a | b} = 0 \\
    \text{Since } a \neq b, \text{ it means } \braket{a | b} = 0.
\end{lefteqs}

The uncoupled basis states are normalized, so the inner product is 1 when the states are identical.

Therefore, in our case:
\begin{lefteqs}
    \bra{j_1 \, j_2 \, m_1 \, m_2} \sum_{m_1', m_2'} C^{j_1 j_2 J}_{m_1' m_2' M} \ket{j_1 \, j_2 \, m_1' \, m_2'} = \\
    \sum_{m_1', m_2'} C^{j_1 j_2 J}_{m_1' m_2' M} \braket{j_1 \, j_2 \, m_1 \, m_2 | j_1 \, j_2 \, m_1' \, m_2'} \\
    = \sum_{m_1', m_2'} C^{j_1 j_2 J}_{m_1' m_2' M} \delta_{m_1 m_1'} \delta_{m_2 m_2'} = C^{j_1 j_2 J}_{m_1 m_2 M}
\end{lefteqs}

And so we've shown that:

\begin{lefteqs}
    \braket{j_1 \, j_2 \, m_1 \, m_2 | J M} = C^{j_1 j_2 J}_{m_1 m_2 M}
\end{lefteqs}

\hrule
\vspace{2em}

\textbf{(b)}
$A^m_l$ and $B^m_l$ are defined in eq. 4.121 (and 4.120) as $A^m_l = \hbar \sqrt{(l - m)(l + m + 1)}$
and $B^m_l = \hbar \sqrt{(l + m)(l - m + 1)}$

Applying $\hat{J}_{\pm} = \hat{J}^{(1)}_{\pm} + \hat{J}^{(2)}_{\pm}$ to equation 6.64 we get:

\begin{lefteqs}
    \hat{J}_{\pm} \ket{J M} = \sum_{m_1, m_2} C^{j_1 j_2 J}_{m_1, m_2, M} \left[\hat{J}^{(1)}_{\pm} \ket{j_1 \, j_2 \, m_1 \, m_2} + \hat{J}^{(2)}_{\pm} \ket{j_1 \, j_2 \, m_1 \, m_2} \right]
\end{lefteqs}

Taking the upper sign ($\hat{J}_+$):

\begin{lefteqs}
    A^M_J \ket{J \, M+1} = \sum_{m_1, m_2} C^{j_1 j_2 J}_{m_1, m_2, M}
    \left[A^{m_1}_{j_1} \ket{j_1 \, j_2 \, m_1+1 \, m_2} + A^{m_2}_{j_2} \ket{j_1 \, j_2 \, m_1 \, m_2+1}  \right]
\end{lefteqs}

Like in part (a), we multiply from the left by $\bra{j_1 \, j_2 \, m_1' \, m_2'}$:

\begin{lefteqs}
    \bra{j_1 \, j_2 \, m_1' \, m_2'} A^M_J \ket{J \, M+1} = \\
    \bra{j_1 \, j_2 \, m_1' \, m_2'} \sum_{m_1, m_2} C^{j_1 j_2 J}_{m_1, m_2, M}
    \left[A^{m_1}_{j_1} \ket{j_1 \, j_2 \, m_1+1 \, m_2} + A^{m_2}_{j_2} \ket{j_1 \, j_2 \, m_1 \, m_2+1}  \right]
\end{lefteqs}

$A^M_J$ is a number, so by part (a) the LHS is
$A^M_J \braket{j_1 \, j_2 \, m_1' \, m_2' | J \, M+1} = A^M_J \, C^{j_1 j_2 J}_{m_1', m_2', M+1}$.
For the RHS, orthonormality of uncoupled eigenstates gives:

\begin{lefteqs}
    A^M_J \, C^{j_1 j_2 J}_{m_1', m_2', M+1} =
    \sum_{m_1, m_2} C^{j_1 j_2 J}_{m_1, m_2, M}
    \left[ A^{m_1}_{j_1} \delta_{m_1 + 1, m_1'} \delta_{m_2, m_2'} + A^{m_2}_{j_2} \delta_{m_1, m_1'} \delta_{m_2 + 1, m_2'} \right] \\
    = A^{m_1'-1}_{j_1} C^{j_1 j_2 J}_{m_1'-1, m_2', M} + A^{m_2'-1}_{j_2} C^{j_1 j_2 J}_{m_1', m_2'-1, M}
\end{lefteqs}

And since $A^{m-1}_l = B^m_l$, this means that:

\begin{lefteqs}
    A^M_J C^{j_1 j_2 J}_{m_1, m_2, M+1} = B^{m_1}_{j_1} C^{j_1 j_2 J}_{m_1 - 1, m_2, M} + B^{m_2}_{j_2} C^{j_1 j_2 J}_{m_1, m_2 - 1, M}
\end{lefteqs}

This is the first recursion relation. The second one is produced in a similar manner, taking the lower sign ($\hat{J}_-$):

\begin{lefteqs}
    B^M_J \ket{J \, M-1} = \sum_{m_1, m_2} C^{j_1 j_2 J}_{m_1, m_2, M}
    \left[B^{m_1}_{j_1} \ket{j_1 \, j_2 \, m_1-1 \, m_2} + B^{m_2}_{j_2} \ket{j_1 \, j_2 \, m_1 \, m_2-1}  \right]
\end{lefteqs}

Multiplying from the left by $\bra{j_1 \, j_2 \, m_1' \, m_2'}$ and using part (a) on the LHS as before:

\begin{lefteqs}
    B^M_J \, C^{j_1 j_2 J}_{m_1', m_2', M-1} =
    \sum_{m_1, m_2} C^{j_1 j_2 J}_{m_1, m_2, M} \left[B^{m_1}_{j_1} \delta_{m_1-1, m_1'} \delta_{m_2, m_2'} + B^{m_2}_{j_2} \delta_{m_1, m_1'} \delta_{m_2-1, m_2'} \right] \\
    = B^{m_1'+1}_{j_1} C^{j_1 j_2 J}_{m_1'+1, m_2', M} + B^{m_2'+1}_{j_2} C^{j_1 j_2 J}_{m_1', m_2'+1, M} \\
    B^M_J C^{j_1 j_2 J}_{m_1, m_2, M-1} = B^{m_1+1}_{j_1} C^{j_1 j_2 J}_{m_1+1, m_2, M} + B^{m_2+1}_{j_2} C^{j_1 j_2 J}_{m_1, m_2+1, M}
\end{lefteqs}

And again, since $B^{m+1}_l = A^m_l$, we can write the relation as:

\begin{lefteqs}
    B^M_J C^{j_1 j_2 J}_{m_1, m_2, M-1} = A^{m_1}_{j_1} C^{j_1 j_2 J}_{m_1+1, m_2, M} + A^{m_2}_{j_2} C^{j_1 j_2 J}_{m_1, m_2+1, M}
\end{lefteqs}


\newpage


\end{document}